% =============================================================================
%  2-Desarrollo.tex — Bitácora de avances y autoevaluación
%
%  INSTRUCCIONES GENERALES:
%  ─────────────────────────────────────────────────────────────────────────────
%  Este archivo contiene dos tablas pre-configuradas:
%
%  TABLA 1 (longtable) — Bitácora cronológica de commits/tags
%    • Llena una fila por cada hito significativo del git log.
%    • Agrupa las filas por semana o etapa usando las filas de encabezado
%      de color (multicolumn). Añade o elimina semanas según necesites.
%    • Obtén los hashes y fechas con:
%        cd <TESIS_DIR> && git log --oneline --format="%h %ad %s" --date=short
%
%  TABLA 2 (tabularx) — Autoevaluación por capítulo o módulo
%    • Ajusta las filas a los capítulos/módulos de tu repositorio.
%    • Las columnas son: Capítulo | Avance concreto | Dificultad | Estado actual
%
%  Parámetros útiles disponibles:
%    \NombreRepo, \VerRef, \VerActual  (definidos en main.tex)
% =============================================================================

\section{Desarrollo y Avances}
\label{sec:desarrollo}

A través del sistema de control de versiones del repositorio \texttt{\NombreRepo},
es posible rastrear con precisión la evolución del proyecto entre las versiones
\texttt{\VerRef} y \texttt{\VerActual}. El Cuadro~\ref{tab:bitacora_git}
presenta una bitácora cronológica de los hitos registrados.

% ─────────────────────────────────────────────────────────────────────────────
%  TABLA 1: Bitácora cronológica (git log)
%
%  Anchos de columna: suma + 2×\tabcolsep×4cols debe ser ≤ \textwidth (155mm)
%  Con \tabcolsep=6pt: 4cols×12pt = 48pt = 16.9mm → columnas ≤ 138mm = 13.8cm
%  Distribución actual: 2.4+2.0+3.3+5.8 = 13.5cm → total 13.5+1.69 = 15.19cm ✓
% ─────────────────────────────────────────────────────────────────────────────
\begin{longtable}{>{\bfseries}p{2.4cm} p{2.0cm} p{3.3cm} p{5.8cm}}
    \caption{Bitácora cronológica de avances en \texttt{\NombreRepo}
             (Corte: \FechaEntrega).}
    \label{tab:bitacora_git} \\
    \toprule
    \textbf{Versión / Hito} & \textbf{Fecha} & \textbf{Ámbito} & \textbf{Descripción del avance} \\
    \midrule
    \endfirsthead

    \multicolumn{4}{l}{\footnotesize\textit{(Continuación del Cuadro~\ref{tab:bitacora_git})}} \\
    \toprule
    \textbf{Versión / Hito} & \textbf{Fecha} & \textbf{Ámbito} & \textbf{Descripción del avance} \\
    \midrule
    \endhead

    \midrule
    \multicolumn{4}{r}{\footnotesize\textit{Continúa en la siguiente página}} \\
    \endfoot

    \bottomrule
    \endlastfoot

    % ── Etapa 1 ───────────────────────────────────────────────────────────────
    \multicolumn{4}{l}{\small\color{ColorPrincipal}\textit{Etapa 1 — Estructura inicial y reencuadre de objetivos (14--22 de abril de 2026)}} \\
    \addlinespace[2pt]

    Config. base
    & 14 Abr 2026
    & Repositorio / CI
    & Configuración de Git LFS para archivos de datos y notebooks; corrección de URLs de badges e incorporación de plantillas de Issues (\texttt{2f9b431}). \\
    \addlinespace[3pt]

    \texttt{v0.3.0} — Reencuadre
    & 14 Abr 2026
    & Cap.\ 1 Introducción
    & \textbf{Hito}: recalibración del alcance a la ZMVM y reescritura de objetivos; integración de cambios de la rama de desarrollo (\texttt{071ecc1}). \\
    \addlinespace[3pt]

    Grafos matemáticos
    & 22 Abr 2026
    & Cap.\ 3.2 Grafos
    & Capítulo de Grafos Matemáticos completado (§§\,3.2.0–3.2.6): teoría de Euler, grafos dirigidos e implementación computacional (\texttt{4de0d6a}). \\
    \addlinespace[8pt]

    % ── Etapa 2 ───────────────────────────────────────────────────────────────
    \multicolumn{4}{l}{\small\color{ColorPrincipal}\textit{Etapa 2 — Capas de Código y Resultados (09--18 de mayo de 2026)}} \\
    \addlinespace[2pt]

    Apimetro v1
    & 09 May 2026
    & Cap.\ 4.1 Apimetro
    & Primera versión del capítulo Apimetro en Capas de Código; descripción de arquitectura y stack tecnológico (\texttt{fe36c71}). \\
    \addlinespace[3pt]

    Cap.\ 4 y 5 --- draft
    & 15 May 2026
    & Caps.\ 4.2 y 5
    & Redacción inicial de VFTModel, estructuración de Cap.\ 5 Resultados e incorporación de imágenes de resultados (\texttt{b252695}, \texttt{89fa819}, \texttt{b01430e}). \\
    \addlinespace[3pt]

    Resultados Fase 1
    & 17 May 2026
    & Cap.\ 5 Fase 1
    & \textbf{Hito}: resultados de Cobertura y Factor de Desviación terminados; integración de Caps.\ 4 y 5 en rama principal (\texttt{394a470}, \texttt{2472d14}). \\
    \addlinespace[3pt]

    LFS y CI
    & 18 May 2026
    & Repositorio / CI
    & Conversión de imágenes de Caps.\ 3 y 4 a Git LFS; habilitación de LFS en pipeline de CI (\texttt{5ede766}). \\
    \addlinespace[8pt]

    % ── Etapa 3 ───────────────────────────────────────────────────────────────
    \multicolumn{4}{l}{\small\color{ColorPrincipal}\textit{Etapa 3 — Ajustes post-Coloquio 2026-2 (20 de agosto -- 02 de septiembre de 2026)}} \\
    \addlinespace[2pt]

    Mapas y Anexo E
    & 20 Ago 2026
    & Cap.\ 5 / Anexos
    & Primer conjunto de ajustes post-coloquio: integración de mapas cartográficos en Cap.\ 5, condensación de secciones de sistemas marginales y creación de Anexo E (\texttt{45d6f55}). \\
    \addlinespace[3pt]

    Mapas integrados
    & 26 Ago 2026
    & Cap.\ 5 Fase 1
    & Integración definitiva de 4 mapas cartográficos en resultados de Fase 1; condensación de 3 bloques técnicos con referencia a Anexo E (\texttt{1a1137d}). \\
    \addlinespace[3pt]

    Correcciones Coloquio
    & 01 Sep 2026
    & Caps.\ 1--5
    & \textbf{Hito}: 36 correcciones ortográficas y de redacción (J-01–J-37) del Dr.\ Jairo aplicadas en 11 archivos; eliminación de bloques \texttt{\textbackslash notaRevision}; compilación limpia a 162~pp.\ (\texttt{a1bf47e}). \\
    \addlinespace[3pt]

    TiKz\_Libraries
    & 02 Sep 2026
    & Caps.\ 4--5 / Figuras
    & Creación de la librería vectorial \texttt{Figures/TiKz\_Libraries/} con 8 diagramas TiKZ migrados y nuevos (arquitectura hexagonal, ports-adapters, pipeline VFT, árbol de directorios, entre otros) (\texttt{b9584ad}). \\
    \addlinespace[3pt]

    Notas de revisión
    & 02 Sep 2026
    & Documentación interna
    & Registro formal de las notas del Dr.\ Jairo (J-01–J-37) y de las 9 observaciones estratégicas de la consulta del 27-Ago (E-01–E-09) en \texttt{Anexos/Notas\_Correciones/} (\texttt{76fbbc8}). \\

\end{longtable}

% ─────────────────────────────────────────────────────────────────────────────
%  TABLA 2: Autoevaluación por capítulo / módulo
% ─────────────────────────────────────────────────────────────────────────────
\vspace{1em}

El Cuadro~\ref{tab:autoevaluacion} presenta una autoevaluación del estado actual
por capítulo o módulo, identificando avances, dificultades encontradas y
los apartados que requieren atención en la siguiente etapa.

\begin{table}[H]
    \centering
    \caption{Autoevaluación por capítulo: avances, redefiniciones y apartados débiles.}
    \label{tab:autoevaluacion}
    \fontsize{9pt}{11.5pt}\selectfont
    \begin{tabularx}{\textwidth}{
        >{\hsize=0.55\hsize\raggedright\arraybackslash}X
        >{\hsize=1.15\hsize\raggedright\arraybackslash}X
        >{\hsize=1.15\hsize\raggedright\arraybackslash}X
        >{\hsize=1.15\hsize\raggedright\arraybackslash}X}
        \toprule
        \textbf{Capítulo / Módulo} &
        \textbf{Avance concreto} &
        \textbf{Dificultad / Redefinición} &
        \textbf{Estado actual y pendientes} \\
        \midrule

        \textbf{Cap.\ 1}\newline Introducción &
        Reencuadre de objetivos al alcance ZMVM; 7 correcciones del Dr.\ Jairo aplicadas (tildes diacríticas, comillas angulares, comas). &
        El objetivo general aún tiene énfasis computacional; requiere reorientación al discurso urbanístico. &
        \textit{Estado:} completo con correcciones.\newline
        \textit{Pendiente:} reformular objetivo y pregunta de investigación (Issue \#12); párrafo de cierre de Fase~2 en §Alcances (Issue \#11). \\
        \addlinespace\hline\addlinespace

        \textbf{Cap.\ 2}\newline Marco Teórico &
        Correcciones de redacción del Dr.\ Jairo aplicadas (\textit{dependencia del automóvil}, \textit{ésta}→\textit{esta}). &
        El marco teórico carece de referentes urbanísticos críticos señalados en el coloquio. &
        \textit{Estado:} completo a nivel descriptivo.\newline
        \textit{Pendiente:} incorporar Lefebvre y Harvey con 2--3 párrafos y citas formales (Issue \#16). \\
        \addlinespace\hline\addlinespace

        \textbf{Cap.\ 3}\newline Conceptos e Indicadores &
        Cap.\ 3.2 (Grafos Matemáticos) terminado en su totalidad. Correcciones J-01, J-13–J-17 aplicadas en los tres subcapítulos. &
        Falta definir explícitamente las capas de Medio Físico Natural señaladas por la evaluadora. &
        \textit{Estado:} sólido; sin bloques pendientes de datos.\newline
        \textit{Pendiente:} nueva subsección de capas físicas en 3-1-Definiciones (Issue \#15). \\
        \addlinespace\hline\addlinespace

        \textbf{Cap.\ 4}\newline Capas de Código &
        Apimetro y VFTModel terminados; 13 correcciones del Dr.\ Jairo aplicadas; diagramas migrados a \texttt{TiKz\_Libraries}; código técnico movido a Anexos A y B. &
        La sección VFTModel no cierra con aplicación en planeación urbana, lo que debilita el argumento ante el coloquio. &
        \textit{Estado:} avanzado.\newline
        \textit{Pendiente:} párrafo ``Aplicación en Planeación Urbana'' al final de §VFTModel (Issue \#13). \\
        \addlinespace\hline\addlinespace

        \textbf{Cap.\ 5}\newline Resultados &
        Fase 1 completa con 4 mapas cartográficos integrados; correcciones J-29–J-37 aplicadas; bloques técnicos condensados con referencia a Anexo E; diagramas de pipeline migrados a librería TiKZ. &
        Fases 2 y 3 permanecen bajo la marca \texttt{\textbackslash bajoRevision} por indicadores pendientes de cálculo (Node Score, T, B, $\Delta E$). &
        \textit{Estado:} Fase 1 \checkmark; Fases 2 y 3 incompletas.\newline
        \textit{Pendiente:} datos de Node Score para Fase 2; protocolo de simulación para Fase 3; crear Anexo D (Issue \#17). \\
        \addlinespace\hline\addlinespace

        \textbf{Cap.\ 6}\newline Análisis Socioespacial &
        Archivo creado en la estructura del documento. &
        Capítulo completamente vacío; es el nexo entre los resultados técnicos y el discurso urbanístico que el coloquio señaló como ausente. &
        \textit{Estado:} vacío --- solo directiva root.\newline
        \textit{Pendiente:} redacción completa una vez acordada la estructura con el asesor (Issue \#8). \\
        \addlinespace\hline\addlinespace

        \textbf{Cap.\ 7}\newline Conclusiones &
        Archivo creado en la estructura del documento. &
        No se puede redactar hasta que Cap.\ 6 y los pendientes de Fases 2/3 estén cerrados. &
        \textit{Estado:} vacío --- bloqueado por dependencias.\newline
        \textit{Pendiente:} redacción final (última etapa). \\
        \bottomrule
    \end{tabularx}
\end{table}

\vspace{1em}
\textit{Nota metodológica:} Para respaldar esta autoevaluación, en la sección
siguiente se presenta evidencia generada con \texttt{latexdiff}, donde las
adiciones aparecen en \textcolor{blue}{azul} y las eliminaciones en
\textcolor{red}{rojo}.
